\subsection{Dynkin graph}

ここでは$C \in \M_n(\Z),\ I = \{1, 2, \ldots, n\}$とする．

\begin{definition}
  \label{def:dynkinGraph}
  \lean{dynkinGraph}
  \leanok
  $C$はGCMであるとする．
  頂点が$V := I$である無向単純グラフを考える(集合としては同じだが，添え字集合$I$を頂点集合として捉えている)．
  $C_{ij} \ne 0$かつ$i \ne j$ならば，頂点$i$と頂点$j$を辺で結ぶ．
  このグラフを$C$の\textbf{Dynkin graph}という．
  GCMの定義の$C_{ij} = 0 \iff C_{ji} = 0$より無向性が従う．
  また，辺があるのは$i \ne j$であるときなので，自己ループがないことも従う．
\end{definition}

\begin{lemma}
  \label{lem:dynkinGraph_preconnected}
  \lean{dynkinGraph_preconnected}
  \leanok
  \uses{}
  $C$はGCMでindecomposableであるとする．
  このとき，$C$のDynkin graphは連結である．
\end{lemma}

\begin{proof}
  ある頂点$u, v \in V$があり，$u$から$v$への道がないとする．
  $S$は$u$から道がある頂点全体の集合とする．
  $u$から$u$へは道があるため，$u \in S$であり，$S$は空集合でない．
  
  $S = V$であることを示す．
  $w \in V$で$w \notin S$なるものがあるとする．
  すると，$S \ne V$である．
  $S$は空集合でも全体集合でもないため，indecomposableの定義からある$i \in S$と$j \notin S$が存在して$C_{ij} \ne 0$である．
  このとき，$i \ne j$であるから，$i$と$j$は辺で結ばれている．
  すると，$u$から$i$への道と$i$と$j$を結ぶ辺から，$u$から$j$への道がある．
  しかし，これは$j \in S$であることを意味し，矛盾する．

  したがって，$S = V$であるから無論$v \in S$である．
  しかし，これは$u$から$v$への道があることを意味し，矛盾する．
\end{proof}

\begin{definition}
  \label{def:neighborSet}
  \lean{neighborSet}
  \leanok
  集合$N(i)$を次で定める：
  \begin{equation}
    N(i) = N_C(i) := \set{j}{C_{ij} \ne 0\ \land\ i \ne j}.
  \end{equation}
  つまり，$i$に隣接している頂点全体の集合である．
  グラフ理論の言葉で\textbf{近傍}という．
  今後，$i$と$j$が隣接しているとき，$i \sim j$と書くこともある．
\end{definition}

\begin{definition}
  \label{def:degree}
  \lean{degree}
  \leanok
  $\deg(i) \in \N$を次のように定め，$C$の頂点$i$の\textbf{degree}という：
  \begin{equation}
    \deg_C(i) = \deg(i) := \# N(i).
  \end{equation}
  つまり，$C$のDynkin graphにおいて頂点$i$に隣接している頂点の数である．
\end{definition}

\begin{definition}
  \label{def:numOfBranch}
  \lean{numOfBranch}
  \leanok
  $n(C) \in \N$を次のように定め，$C$の\textbf{branchの数}という：
  \begin{equation}
    n(C) := \# \set{i}{3 \le \deg(i)}.
  \end{equation}
  つまり，$C$のDynkin graphにおいてY字や十字以上に分岐している頂点の数である．
\end{definition}